\paragraph{Adaptive no-gain and sufficient statistics: sampling without replacement and counting posterior under microcanonical prior}
In the following, under the uniform microcanonical prior \(F\sim\mathrm{Unif}(\mathfrak{Fold}_m(d))\),
we rigorously equate any deterministic adaptive query process to sampling without replacement from the multiset \(\{x^{\times d(x)}\}_{x\in X_m}\);
and show that the counting vector \(C_k\) constitutes a minimal sufficient statistic for \(F\) in the posterior sense.

\begin{definition}[Adaptive distinct query trajectory and counting statistics]\label{def:pom-microcanonical-adaptive-trajectory-counts}
Let \(\mathcal{A}\) be any deterministic adaptive algorithm performing \(k\) distinct queries: at step \(i\) it selects
\[
\omega_i\in\Omega_m\setminus\{\omega_1,\dots,\omega_{i-1}\},
\]
and observes the response
\[
T_i:=F(\omega_i)\in X_m.
\]
Denote the trajectory \(T_{1:k}:=(T_1,\dots,T_k)\in X_m^k\), and define the counting vector
\[
C_k(x):=\card{\{i\in\{1,\dots,k\}:\ T_i=x\}},\qquad x\in X_m.
\]
\end{definition}

\begin{theorem}[Adaptive no-gain: trajectory distribution and stepwise sampling-without-replacement law]\label{thm:pom-microcanonical-adaptive-no-gain-without-replacement}
In the setting of Definition \ref{def:pom-microcanonical-adaptive-trajectory-counts}:
\begin{enumerate}
  \item \textbf{(Trajectory distribution independent of strategy)}
  For any \(t=(t_1,\dots,t_k)\in X_m^k\), its occurrence probability depends only on \(k\) and the fiber spectrum \(d\), and satisfies
  \[
  \boxed{\;
  \mathbb{P}_{\mathcal{A}}(T_{1:k}=t)
  =\frac{1}{(N)_k}\prod_{x\in X_m}\bigl(d(x)\bigr)_{c_t(x)}\;},
  \]
  where \(c_t(x)=\card{\{i:\ t_i=x\}}\) is the counting vector of \(t\), \((u)_r:=u(u-1)\cdots(u-r+1)\) is the falling factorial,
  and we adopt the convention \((N)_k:=N(N-1)\cdots(N-k+1)\).
  \item \textbf{(Stepwise sampling without replacement)}
  For each \(i\ge 1\) and any \(x\in X_m\), the conditional probability formula holds
  \[
  \boxed{\;
  \mathbb{P}_{\mathcal{A}}\!\bigl(T_i=x\mid T_{1:i-1}\bigr)
  =\frac{d(x)-C_{i-1}(x)}{N-i+1}\; } .
  \]
  Hence \((T_i)\) is identically distributed to the sampling-without-replacement process from the multiset \(\{x^{\times d(x)}\}_{x\in X_m}\).
\end{enumerate}
\end{theorem}
\begin{proof}
For any given trajectory \(t\in X_m^k\), the query sequence \((\omega_i)\) of algorithm \(\mathcal{A}\) is a deterministic function of the prefix of \(t\);
therefore the event \(\{T_{1:k}=t\}\) is equivalent to certain distinct point set \(Q_t=\{\omega_1(t),\dots,\omega_k(t)\}\subset\Omega_m\) satisfying
\(
F(\omega_i(t))=t_i
\).
By the trajectory probability formula of Theorem \ref{thm:pom-microcanonical-fold-posterior-count-and-prob},
\(\mathbb{P}(T_{1:k}=t)\) depends only on \(c_t\) and gives the shown expression; this expression contains no geometric information about \(Q_t\), thus is also independent of the choice rule of \(\mathcal{A}\).
\par
The stepwise conditional probability is obtained by ratio calculation: for any history prefix \(t_{1:i-1}\) and any \(x\in X_m\),
applying the probability formula of Theorem \ref{thm:pom-microcanonical-fold-posterior-count-and-prob} respectively to the length-\(i\) trajectory \((t_{1:i-1},x)\) and the length-\((i-1)\) trajectory \(t_{1:i-1}\), and taking the ratio gives
\[
\mathbb{P}(T_i=x\mid T_{1:i-1}=t_{1:i-1})
=\frac{(d(x))_{c_{t_{1:i-1}}(x)+1}}{(d(x))_{c_{t_{1:i-1}}(x)}}\cdot\frac{(N)_{i-1}}{(N)_{i}}
=\frac{d(x)-C_{i-1}(x)}{N-i+1},
\]
where we used the identity \((u)_{r+1}=(u-r)(u)_r\).
The conclusion holds. \qed
\end{proof}

\begin{theorem}[Counting sufficient statistic and posterior uniformity]\label{thm:pom-microcanonical-count-sufficient-statistic-posterior-uniform}
In the setting of Definition \ref{def:pom-microcanonical-adaptive-trajectory-counts}, suppose the algorithm obtains trajectory \(t\in X_m^k\) after \(k\) queries, and let \(Q_k:=\{\omega_1,\dots,\omega_k\}\),
\(U_k:=\Omega_m\setminus Q_k\).
\begin{enumerate}
  \item \textbf{(Posterior depends only on counts)}
  For any two trajectories \(t,t'\in X_m^k\), if \(c_t=c_{t'}\), then the induced posterior distributions are isomorphic:
  there exists a domain permutation \(\sigma\in\mathrm{Sym}(\Omega_m)\) such that
  \[
  \mathrm{Law}\bigl(F\mid T_{1:k}=t\bigr)\ \cong\ \mathrm{Law}\bigl(F\circ\sigma^{-1}\mid T_{1:k}=t'\bigr).
  \]
  Hence all influence of the trajectory on the posterior can be compressed to the counting vector \(C_k\), and \(C_k\) is a minimal sufficient statistic for \(F\) in this sense.
  \item \textbf{(Conditional posterior uniformity on unqueried part)}
  Given trajectory \(t\), the posterior restricted to the unqueried set \(U_k\) is uniform:
  \[
  \mathrm{Law}\bigl(F|_{U_k}\mid T_{1:k}=t\bigr)
  =\mathrm{Unif}\bigl(\mathfrak{F}_{U_k}(d-c_t)\bigr),
  \]
  where \(\mathfrak{F}_{U_k}(d-c_t)\) denotes the set of all mappings from \(U_k\) to \(X_m\) with fiber sizes \(d(x)-c_t(x)\).
\end{enumerate}
\end{theorem}
\begin{proof}
For any given query set \(Q_k\) and label assignment \(t\) on it, the compatible set
\[
\Bigl\{F\in\mathfrak{Fold}_m(d):\ F(\omega_i)=t_i,\ 1\le i\le k\Bigr\}
\]
only requires realizing the residual count \(d-c_t\) on \(U_k\);
since the prior is uniform on \(\mathfrak{Fold}_m(d)\), after conditioning it remains uniform on this compatible set.
Restricting this uniformity to \(U_k\) yields item (2).
\par
If two trajectories \(t,t'\) have the same counting vector, there exists a permutation that transports the label assignment on \(Q_k\) from \(t\) to \(t'\) without changing the usage count of each label;
by the permutation invariance of the prior on \(\Omega_m\) we can identify the two posteriors as isomorphic, yielding item (1).
Minimal sufficiency follows directly from "posterior depends only on \(c_t\)" and "residual count \(d-c_t\) uniquely determines the posterior feasible set." \qed
\end{proof}

\endinput

